\chapter{Параллельные вычисления в задаче моделирования теплопередачи}

\section{Математическая модель нестационарной теплопроводности в трёхмерной области}

\subsection{Уравнение теплопроводности с граничными условиями первого рода}

Рассматривается однородное изотропное твёрдое тело в форме куба $\Omega = [0, L]^3$ с длиной ребра $L = 1$~м. Распределение температуры $T(x, y, z, t)$ в нём при отсутствии внутренних источников тепла описывается уравнением теплопроводности Фурье--Кирхгофа~\cite{сергиенко2004исследование}
\begin{equation}
    \frac{\partial T}{\partial t} = a \left( \frac{\partial^2 T}{\partial x^2} + \frac{\partial^2 T}{\partial y^2} + \frac{\partial^2 T}{\partial z^2} \right),
    \label{eq:heat}
\end{equation}
где $a$ --- коэффициент температуропроводности материала, м$^2$/с. Этот коэффициент равен отношению теплопроводности к произведению плотности на удельную теплоёмкость и характеризует скорость выравнивания температуры в теле~\cite{vlasov2004research}. Для конструкционной стали, которая используется в расчётах данной работы, принято значение $a = 1{,}28 \cdot 10^{-5}$~м$^2$/с.

Для однозначного решения уравнения~\eqref{eq:heat} задаются начальное и граничные условия~\cite{luo2010effect}. В начальный момент температура во всех точках тела одинакова и равна $T_0$. На каждой из шести граней куба поддерживается собственная постоянная температура, то есть используются граничные условия первого рода, или условия Дирихле~\cite{дирихлерешении}. Например, на гранях, перпендикулярных оси $x$, они имеют вид
\begin{equation}
    T\big|_{x=0} = T_{x-}, \qquad T\big|_{x=L} = T_{x+},
\end{equation}
а на остальных гранях записываются аналогично. Независимое задание температуры на каждой грани позволяет моделировать как равномерный прогрев тела, так и направленный тепловой поток. Подробный вывод уравнения и обсуждение физического смысла его параметров приведены в отчёте за предыдущий семестр, поэтому здесь математическая модель изложена кратко.

\subsection{Явная конечно-разностная схема и условие её устойчивости}

Расчётная область покрывается равномерной сеткой из $N_x \times N_y \times N_z$ узлов с шагами $h_x = L / (N_x - 1)$, $h_y$ и $h_z$, которые определяются аналогично. Обозначим через $T^n_{i,j,k}$ приближённое значение температуры в узле с индексами $(i, j, k)$ в момент времени $t_n = n \Delta t$. Замена производной по времени разностью вперёд, а вторых производных по пространству центральными разностями приводит к явной схеме FTCS~\cite{kolgan1975finite}
\begin{equation}
    T^{n+1}_{i,j,k} = T^n_{i,j,k} + r_x \, \delta_x^2 T^n_{i,j,k} + r_y \, \delta_y^2 T^n_{i,j,k} + r_z \, \delta_z^2 T^n_{i,j,k},
    \label{eq:ftcs}
\end{equation}
в которой использованы обозначения $\delta_x^2 T_{i,j,k} = T_{i+1,j,k} - 2T_{i,j,k} + T_{i-1,j,k}$ и $r_x = a \Delta t / h_x^2$, а величины для осей $y$ и $z$ определяются так же. Схема~\eqref{eq:ftcs} имеет первый порядок аппроксимации по времени и второй порядок по пространству.

Явная схема устойчива только при ограничении на шаг по времени~\cite{dubey2015stability}. Для трёхмерной задачи оно записывается в виде неравенства
\begin{equation}
    2 \left( r_x + r_y + r_z \right) \le 1,
    \label{eq:stab}
\end{equation}
поэтому в программе шаг $\Delta t$ выбирается автоматически так, чтобы левая часть~\eqref{eq:stab} равнялась 0,4. Такой запас исключает рост ошибок округления при любых допустимых параметрах расчёта. Из неравенства~\eqref{eq:stab} следует, что шаг по времени пропорционален квадрату шага сетки, и именно это делает расчёт на мелких сетках трудоёмким.

Для дальнейшего важно одно свойство схемы~\eqref{eq:ftcs}. Новое значение температуры в узле вычисляется только по значениям с предыдущего временного слоя в самом узле и в шести соседних узлах. Следовательно, на одном шаге по времени все внутренние узлы можно обрабатывать независимо друг от друга и в произвольном порядке, что и создаёт возможность для параллельных вычислений.

\section{Технология OpenMP для вычислительных систем с общей памятью}

\subsection{Модель исполнения fork--join и директивы распараллеливания циклов}

OpenMP представляет собой стандарт программного интерфейса для параллельного программирования на языках C, C++ и Fortran в системах с общей памятью~\cite{dagum1998openmp}. Стандарт включает директивы компилятора, библиотечные функции и переменные окружения\footnote{Актуальные версии спецификации OpenMP публикуются на официальном сайте OpenMP Architecture Review Board: \url{https://www.openmp.org}.}. Программа начинает работу в одном главном потоке, а при входе в параллельную область создаётся группа потоков, которые выполняют её совместно. По завершении области потоки синхронизируются, и выполнение продолжается в главном потоке. Такая схема называется моделью fork--join~\cite{chandra2001parallel}.

В вычислительных программах основная работа сосредоточена в циклах, поэтому наиболее употребительна директива \texttt{parallel for}. Она распределяет итерации следующего за ней цикла между потоками группы, причём программисту не требуется вручную создавать потоки и делить диапазон индексов~\cite{antonov2009openmp}. Все потоки работают в общем адресном пространстве, поэтому массивы доступны каждому из них без копирования. Обратной стороной общей памяти является риск гонки данных, при которой несколько потоков одновременно изменяют одну и ту же ячейку и результат зависит от порядка их выполнения. Если программа компилируется без поддержки OpenMP, директивы игнорируются и получается обычная последовательная программа, что удобно для сравнения двух версий.

\subsection{Распределение итераций между потоками и редукция данных}

Способ распределения итераций задаётся параметром \texttt{schedule}. При статическом распределении диапазон итераций заранее делится на равные непрерывные части по числу потоков, а при динамическом распределении потоки получают новые порции итераций по мере завершения предыдущих~\cite{antonov2009openmp}. Статическое распределение не требует дополнительной синхронизации во время работы цикла и поэтому предпочтительно, когда все итерации выполняют одинаковый объём вычислений. Для вложенных циклов стандарт, начиная с версии 3.0, предоставляет параметр \texttt{collapse}, объединяющий несколько циклов в одно пространство итераций.

Отдельного внимания требуют циклы, в которых все итерации накапливают результат в одной переменной, например при суммировании элементов массива. Параметр \texttt{reduction} создаёт для каждого потока собственную копию такой переменной, а по окончании цикла объединяет частичные результаты~\cite{chandra2001parallel}. Этот приём устраняет гонку данных без использования блокировок. Однако порядок сложения при этом меняется, и для чисел с плавающей точкой результат может отличаться от последовательного в последних значащих разрядах.

\section{Показатели эффективности параллельных программ и их ограничения}

\subsection{Ускорение и эффективность параллельной программы на многоядерном процессоре}

Качество распараллеливания оценивается по времени выполнения одной и той же задачи на разном числе потоков. Если $t_1$ --- время расчёта одним потоком, а $t_p$ --- время расчёта $p$ потоками, то ускорение и эффективность определяются выражениями
\begin{equation}
    S_p = \frac{t_1}{t_p}, \qquad E_p = \frac{S_p}{p}.
    \label{eq:speedup}
\end{equation}
В идеальном случае ускорение равно числу потоков, а эффективность равна единице. На практике эффективность меньше единицы из-за затрат на создание и синхронизацию потоков, неравномерной загрузки и наличия участков программы, которые выполняются последовательно. Значения эффективности выше 0,8 обычно считаются хорошим результатом для программ с общей памятью.

\subsection{Закон Амдала и ограничение пропускной способностью памяти}

Пусть доля $f$ времени работы последовательной программы приходится на участки, которые не распараллеливаются. Тогда по закону Амдала~\cite{amdahl1967validity} ускорение на $p$ потоках не может превысить величину
\begin{equation}
    S_p = \frac{1}{f + (1 - f)/p},
    \label{eq:amdahl}
\end{equation}
и при неограниченном увеличении числа потоков стремится к $1/f$. Например, при $f = 0{,}1$ ускорение на четырёх потоках не превышает 3,1, а на любом числе потоков --- 10. В программе моделирования теплопередачи последовательной частью является прежде всего запись промежуточных результатов в файлы, поскольку она выполняется одним потоком.

Второе ограничение связано с работой памяти. При вычислении нового значения по формуле~\eqref{eq:ftcs} на каждый узел приходится около десяти арифметических операций и обращение к семи значениям массива. Процессор выполняет такие операции быстрее, чем память успевает доставлять данные, поэтому скорость расчёта определяется пропускной способностью памяти. Ядра процессора разделяют общую шину памяти и кэш последнего уровня~\cite{кудухов2014особенности}, и с ростом числа потоков каждое ядро получает меньшую долю этой пропускной способности. По этой причине для задач со стенсильными вычислениями эффективность обычно снижается с ростом числа потоков даже при отсутствии последовательных участков.
